\section{Synchronization of the anchor coordinates}
\label{sec:anchor-synchronization}

The anchor block has two purposes: low internal energy forces all of its CRT coordinates to come from one integer, while the remaining assignments have small total Fourier weight.

\subsection{A fixed-ratio reciprocal prime code}

Fix $\delta\in(0,1)$. Let $Y$ be large and let $Q$ be a set of primes in $[\delta Y,Y)$ satisfying
\[
 c_\delta\frac{Y}{\log Y}\leq N:=|Q|\leq C_\delta\frac{Y}{\log Y}.
\]
For an assignment $a=(a_p)_{p\in Q}$ and distinct $p,q\in Q$, let $H_{pq}(a)$ be the unique centred CRT lift with
\[
 H_{pq}\equiv a_p\pmod p,\qquad H_{pq}\equiv a_q\pmod q,\qquad -\frac{pq}{2}<H_{pq}\leq\frac{pq}{2}.
\]
Put
\[
 \mathcal Q(a)=\sum_{p<q}\left(\frac{H_{pq}(a)}{pq}\right)^2,
 \qquad
 \sigma_Q^2=\sum_{p<q}\frac1{p^2q^2}\asymp_\delta\frac1{Y^2\log^2Y}.
\]
Fix $\rho\in(0,1/4)$.

\begin{lemma}[Reciprocal dispersion]
\label{lem:reciprocal-dispersion}
There are constants $s_\delta\geq16$ and $c_D(\delta)>0$ such that, if $q$ is prime in $[\delta Y,Y)$, $\mathcal F$ is a set of $s\geq s_\delta$ primes in $[\delta Y,Y)$ not containing $q$, and $d\not\equiv0\pmod q$, then
\[
 \sum_{p\in\mathcal F}\dist{d p^{-1}/q}^2\geq c_D(\delta)\frac{s^3}{Y^2}.
\]
The constants are uniform when $\delta$ ranges in a compact subset of $(0,1)$.
\end{lemma}
\begin{proof}
An interval of length less than $Y$ contains at most $\nu_\delta:=1+\lceil\delta^{-1}\rceil$ integers in one residue class modulo $q$. Put $\Delta=s/(128\nu_\delta Y)$. If $\dist{dp^{-1}/q}\leq\Delta$, choose an integer $\ell$ such that
\[
 v=dp^{-1}-\ell q,\qquad |v|\leq\Delta q\leq\frac{s}{128\nu_\delta}.
\]
Then $vp\equiv d\pmod q$. For each fixed nonzero $v$, this determines one residue class of $p$ modulo $q$, containing at most $\nu_\delta$ elements of $[\delta Y,Y)$. For $s_\delta$ sufficiently large, fewer than $s/2$ primes satisfy the small-distance condition. The others have distance greater than $\Delta$, proving the claim.
\end{proof}

Call an integer $m$ a \emph{dominant label} if $|m|\leq\delta^2Y^2/8$ and at least $(1-\rho)N$ coordinates satisfy $a_p\equiv m\pmod p$. A dominant label is unique: two dominant classes meet in at least two primes, whose product is at least $\delta^2Y^2$, while the difference of the labels has absolute value at most $\delta^2Y^2/4$.

\begin{lemma}[Cross-label energy]
\label{lem:cross-label-energy}
There is $A_\delta>0$ such that the following holds. Let $C_m,C_{m'}$ be disjoint label classes with $m\neq m'$ and
\[
 |m|,|m'|\leq B_0<\frac{\delta^2Y^2}{16}.
\]
If
\[
 |C_m|\geq\max(s_\delta,A_\delta B_0/Y),\qquad |C_{m'}|\geq2,
\]
then
\[
 \sum_{p\in C_m}\sum_{q\in C_{m'}}\left(\frac{H_{pq}}{pq}\right)^2
 \geq c_E(\delta)\frac{|C_m|^3(|C_{m'}|-1)}{Y^2}.
\]
\end{lemma}
\begin{proof}
Put $d=m'-m$. At most one prime in $C_{m'}$ divides $d$. Fix any remaining $q\in C_{m'}$. For $p\in C_m$, write $H_{pq}=m+jp$. Reduction modulo $q$ gives $jp\equiv d\pmod q$ and hence
\[
 \dist{dp^{-1}/q}\leq\frac{|H_{pq}|}{pq}+\frac{|m|}{pq}.
\]
By \cref{lem:reciprocal-dispersion}, at least half of the left sides exceed a constant times $|C_m|/Y$. The size hypothesis makes $|m|/(pq)$ at most half that amount. Thus these $p$ contribute $\gg_\delta |C_m|^3/Y^2$ for each admissible $q$, and summing over $q$ proves the lemma.
\end{proof}

\subsection{Low energy forces one exact label}

\begin{proposition}[Nondominant forcing]
\label{prop:nondominant-forcing}
There is $c_w(\delta)>0$ such that, for all sufficiently large $Y$,
\[
 \mathcal Q(a)<c_w(\delta)\frac{Y}{\log^3Y}
\]
implies that $a$ has a unique dominant label.
\end{proposition}
\begin{proof}
Assume $0<R:=\mathcal Q(a)<cY/\log^3Y$; the case $R=0$ is immediate. Put
\[
 B_0=A_{\delta,\rho}\frac{\sqrt R\,Y^2}{N}
\]
and call a pair $\{p,q\}$ bad if $|H_{pq}|>B_0$. Since $pq<Y^2$, the number of bad pairs is at most $N^2/A_{\delta,\rho}^2$. A base prime $p_0$ therefore has bad degree at most $\rho N/8$ when $A_{\delta,\rho}$ is large.

For every nonbad neighbour $q$, the integer $H_{p_0q}$ is a label for $a_q$, lies in $[-B_0,B_0]$, and is congruent to the fixed residue $a_{p_0}$ modulo $p_0$. Hence at most
\[
 M_0\leq\frac{2B_0}{\delta Y}+2
\]
labels occur. Decreasing $c$ if necessary gives $B_0<\delta^2Y^2/16$.

Suppose no class is dominant. Call a label class small if its size is below
\[
 s_0=C_{\delta,\rho}(B_0/Y+1),
\]
where $C_{\delta,\rho}$ makes $s_0\geq s_\delta$ and $s_0\geq A_\delta B_0/Y$. If small classes contain at least $\rho N/4$ vertices, then
\[
 \frac{\rho N}{4}\leq M_0s_0\leq C_{\delta,\rho}(B_0/Y+1)^2.
\]
Writing $u=B_0/Y$, the case $u<1$ contradicts $N\to\infty$, while $u\geq1$ gives $R\gg_{\delta,\rho}N^3/Y^2\gg_\delta Y/\log^3Y$.

Otherwise remove $p_0$, its bad neighbours, and all small classes. The remaining substantial classes have total size $S\geq(1-\rho/2)N$, each has size at least $s_0$, and no one has size above $(1-\rho)N$. Applying \cref{lem:cross-label-energy} between distinct classes gives
\[
 R\gg_\delta\frac1{Y^2}\sum_i n_i^3\sum_{j\neq i}(n_j-1)
 \gg_{\delta,\rho}\frac{N^4}{M_0^2Y^2}.
\]
If $u<1$, this is $\gg Y^2/\log^4Y$; if $u\geq1$, substituting $M_0^2\ll_{\delta,\rho}RY^2/N^2$ gives $R\gg N^3/Y^2$. Both contradict the assumed upper bound for sufficiently small $c$. Thus a dominant label exists.
\end{proof}

\begin{proposition}[Exact anchor rigidity]
\label{prop:exact-cold-rigidity}
After decreasing $c_w(\delta)$ if necessary, every assignment satisfying
\[
 \mathcal Q(a)<c_w(\delta)\frac{Y}{\log^3Y}
\]
is exactly represented by one integer $m$:
\[
 a_p\equiv m\pmod p\qquad(p\in Q).
\]
Moreover,
\[
 |m|\ll_{\delta,\rho}\frac{\sqrt{\mathcal Q(a)}}{\sigma_Q}.
\]
\end{proposition}
\begin{proof}
Let $C$ be the dominant class with label $m$. Internal pairs in $C$ have centred lift $m$, so
\[
 \mathcal Q(a)\geq m^2\sum_{p<q\in C}\frac1{p^2q^2}\gg_{\delta,\rho}m^2\sigma_Q^2.
\]
This proves the label bound. If $q\notin C$, put $d\equiv a_q-m\pmod q$. As in \cref{lem:cross-label-energy}, reciprocal dispersion against $p\in C$ forces the edges incident with $q$ to contribute $\gg_\delta |C|^3/Y^2$. Distinct exceptional $q$ give disjoint edge sets, whence
\[
 |Q\setminus C|\ll_\delta \mathcal Q(a)Y^2/N^3<C_\delta c_w(\delta).
\]
Choosing $c_w(\delta)$ so that the last quantity is below one leaves no exceptional coordinate.
\end{proof}

\subsection{Fingerprint counting}

\begin{lemma}[Fingerprint rigidity]
\label{lem:fingerprint-rigidity}
Fix $\mathcal F\subset Q$ with $|\mathcal F|=s\geq s_\delta$ and fix its residues. For $q\in Q\setminus\mathcal F$ and $w\pmod q$, put
\[
 t_q(w)=\sum_{p\in\mathcal F}\left(\frac{H_{pq}(a_p,w)}{pq}\right)^2.
\]
For a fixed $g_\delta>0$, at most one residue $w$ satisfies
\[
 t_q(w)<g_\delta\frac{s^3}{Y^2}.
\]
\end{lemma}
\begin{proof}
For distinct $w,w'$ put $d=w-w'$. The difference of the two centred lifts represents $dp^{-1}/q$ on the circle, so
\[
 \dist{dp^{-1}/q}^2\leq2\left(\frac{H_{pq}(a_p,w)}{pq}\right)^2+2\left(\frac{H_{pq}(a_p,w')}{pq}\right)^2.
\]
Sum and use \cref{lem:reciprocal-dispersion}; choosing $4g_\delta<c_D(\delta)$ rules out two small-energy residues.
\end{proof}

\begin{proposition}[Fingerprint counting bound]
\label{prop:fingerprint-counting}
For every $\varepsilon>0$, all sufficiently large fixed-ratio blocks satisfy
\[
 \#\{a:\mathcal Q(a)\leq R\}\leq\exp(\varepsilon R)
\]
uniformly for $R\geq c_w(\delta)Y/\log^3Y$.
\end{proposition}
\begin{proof}
When $R\leq c_{*,\delta}N^4/Y^2$, choose a fixed fingerprint $\mathcal F$ of size
\[
 s=\left\lceil A_\delta(RY^2)^{1/4}\right\rceil.
\]
There are at most $Y^s$ assignments on it. By \cref{lem:fingerprint-rigidity}, all but
\[
 u\leq\frac{RY^2}{g_\delta s^3}\ll_\delta s
\]
remaining vertices are forced. Choosing the exceptions and their residues gives at most
\[
 \exp\bigl(O_\delta((RY^2)^{1/4}\log Y)\bigr)=\exp(o(R))
\]
possibilities in the stated lower range. If $R>c_{*,\delta}N^4/Y^2$, the crude bound $Y^N$ is already $\exp(o(R))$.
\end{proof}

\subsection{Application to the anchor block}

Apply the preceding results with $Y=Z$, $\delta=\eta$, and $Q=B$. For an anchor assignment $a_B$, put
\[
 Q_B(a_B)=\sum_{q<q'\in B}\left(\frac{H_{qq'}(a_B)}{qq'}\right)^2,
 \qquad
 \sigma_{B,0}^2=\sum_{q<q'\in B}\frac1{q^2q'^2}.
\]
The complete-pair identity and \cref{eq:B-power-sums} give
\begin{equation}
\label{eq:anchor-square-scale}
 \sigma_{B,0}^2\sim\frac{(\eta^{-1}-1)^2}{2Z^2\log^2Z}\asymp_\eta\frac1{Z^2\log^2Z}.
\end{equation}
There is a forcing floor
\[
 F_B=c_B(\eta)\frac{Z}{\log^3Z}
\]
such that every anchor assignment with $Q_B<F_B$ has one exact integer label $m$, and then
\begin{equation}
\label{eq:top-exact-energy}
 Q_B(a_B)=m^2\sigma_{B,0}^2.
\end{equation}

Define the anchor modulus
\[
 T_B(a_B)=\prod_{q<q'\in B}\left|(1-\theta)+\theta\exp\left(2\pi i\frac{H_{qq'}(a_B)}{qq'}\right)\right|.
\]
By \cref{eq:bernoulli-modulus}, $T_B(a_B)\leq\exp(-\kappa_{\rm mod}Q_B(a_B))$. The fingerprint bound, summed over unit energy shells, gives
\begin{equation}
\label{eq:noncoherent-top-tail}
 \sum_{Q_B\geq F_B}T_B(a_B)\ll_\eta\exp(-c_\eta F_B).
\end{equation}
Below the floor, distinct assignments have distinct labels, and therefore
\[
 \sum_{Q_B<F_B}T_B(a_B)\leq\sum_{m\in\mathbb Z}\exp(-\kappa_{\rm mod}m^2\sigma_{B,0}^2)\ll_\eta Z\log Z.
\]

\begin{proposition}[Anchor partition]
\label{prop:top-partition}
The complete anchor weight satisfies
\begin{equation}
\label{eq:P-top}
 P_{\mathrm{anc}}:=\sum_{a_B}T_B(a_B)\ll_\eta Z\log Z,
\end{equation}
while its noncoherent part satisfies \cref{eq:noncoherent-top-tail}. Every coherent anchor assignment carries a unique exact integer label and obeys \cref{eq:top-exact-energy}.
\end{proposition}
